\subsection{Subgroup Definitions and Coding for Heterogeneity Analysis}
\label{subsec:subgroup_definitions}

To ensure a unified narrative and comparability across the three essays, the subgroup classifications used for heterogeneity analysis are harmonized based on the common data structure described in Section 2.4. Table~\ref{tab:subgroup_coding} summarizes the primary subgroup definitions applied throughout the dissertation.

\begin{table}[h!]
\centering
\small
\caption{Unified Subgroup Definitions and Analytic Coding}
\label{tab:subgroup_coding}
\begin{tabular}{lp{6cm}c}
\toprule
\textbf{Construct} & \textbf{Operational Definition} & \textbf{Coding Rule} \\
\midrule
Institutional Trust & High: Explicitly expressed confidence (``Strongly trust'' or ``Somewhat trust''). \newline Low: Lack of explicit confidence (``Neutral,'' ``Somewhat distrust,'' or ``Strongly distrust''). & Binary ($1$=High, $0$=Low) \\
\addlinespace
Smoking Status & Smoker: Currently smokes at least occasionally (``Often'' or ``Sometimes''). \newline Non-smoker: Has never smoked (``Never''). & Binary ($1$=Smoker, $0$=Non-smoker) \\
\addlinespace
Physical Activity & Low: Rare activity. \newline Mid: Activity some of the time. \newline High: Activity quite often. & Categorical / Ordinal \\
\addlinespace
Biological Vulnerability & Indicators for chronic conditions and self-rated health scores. & Binary / Categorical \\
\addlinespace
Demographics & Age bins ($\geq$55 vs. $<$55), Education (Primary/Secondary vs. Tertiary), and Urban-Rural residence. & As specified in Section 2.4.3 \\
\bottomrule
\end{tabular}
\end{table}

These consistent definitions support the synthesis of findings in Chapter 6, where behavioral delay sensitivities (Essay 1), monetized delay burdens (Essay 2), and equity-sensitive policy rules (Essay 3) are integrated into a single evaluative framework. Any deviations from these baseline definitions for specific sensitivity checks are explicitly noted in the respective chapters.
